\paragraph{Canonical semantics can require repeated sources.}
Exact dependence information does not make a distinguished input linear.
The obstruction is visible in the Boolean contexts in which that input
matters.

Fix a finite passive variable set $V$ and a distinguished input $x\notin V$.
Write $f(\epsilon,A)$ for the Boolean value with distinguished input
$x=\epsilon$ and passive true set $A\subseteq V$, and put
\[
 \mathcal E_x(f)=\{A\subseteq V:f(0,A)=0,\ f(1,A)=1\}.
\]
Define a graph on $\mathcal E_x(f)$ by joining $P,Q$ when, in either
orientation, some passive cube $[L,U]$ satisfies
\begin{equation}\label{core:eq:essential-context-conflict}
 P\subseteq U,\qquad L\subseteq Q,\qquad
 \mathcal E_x(f)\cap[L,U]=\varnothing.
\end{equation}
On such a cube $x$ is inessential. Its common Boolean cofactor is
nonconstant: monotonicity and the two essential contexts force value
$0$ at $L$ and value $1$ at $U$.

\begin{proposition}[Essential-context coloring]
\label{core:prop:essential-context-coloring}
Let $f$ be a nonconstant monotone Boolean function. Every core term whose
four-state restriction is pointwise contained in $D_f$ uses $x$ at least
as many times as the chromatic number of the graph
\eqref{core:eq:essential-context-conflict}. Choices by additional owners
and ordinary union are allowed.
\end{proposition}
\begin{proof}
Choose distinct owners $r_0,r_1,k$. Use endpoint menus
$\{r_0\},\{r_1\}$, neutral $\nv=\{r_0,r_1\}$, and
marked $\pv=\{r_0,r_1,k\}$. Containment in $D_f$
forces the term to compute $f$ on Boolean inputs. At a context
$E\in\mathcal E_x(f)$, mark $x$ and keep the passive inputs Boolean.
The output contains the marker: otherwise omission independence
(\cref{core:prop:omission}) would
let us replace $x$ by either endpoint without changing the value,
contrary to essentiality. Choose an $x$-occurrence supplying this marker
along a surviving path.

One occurrence cannot serve adjacent contexts $P,Q$. At every
$r_0$-owned gate on its path, the sibling is not the absorbing atom $r_0$
at context $P$. At the crossing cube its $r_0$-atomic input coalition is
$V\setminus U\subseteq V\setminus P$; $x$ is marked in both inputs.
Monotonicity of the sibling's singleton profile keeps that flag false.
At an $r_1$-owned gate use context $Q$ instead: its $r_1$-atomic coalition
shrinks from $Q$ to $L$, so that sibling still is not $r_1$.
All other owners act as union on the four-state subalgebra. Hence no
path gate absorbs the marked source on the crossing cube. The marker
survives, contradicting containment in $D_f$, which there equals $\nv$.
Choosing one surviving occurrence at every essential context therefore
gives a proper coloring by the $x$-occurrences.
\end{proof}

\begin{theorem}[Unbounded source repetition in canonical extensions]
\label{core:thm:canonical-source-repetition}
For $m\geq2$, take disjoint variable pairs $\{a_i,b_i\}$ and a new
input $x$. Put
\[
 A_i=a_i\wedge b_i,\qquad
 O_i=\bigvee_{j\ne i}(a_j\vee b_j),\qquad
 T_m=\bigvee_{i=1}^m\bigl(A_i\wedge(x\vee O_i)\bigr),
\]
and let $F_m$ be its Boolean function. Then $T_m$ realizes $D_{F_m}$,
and every realization of $D_{F_m}$ has at least $m$ occurrences of $x$.
The displayed term attains the bound. This holds for every fixed
$K\geq3$, even when all core choices and union are available.
\end{theorem}
\begin{proof}
Each summand is read-once and therefore canonical. To prove that their
outer join is canonical, consider a nonconstant Boolean cube. No summand
is constantly $1$, so the outer join carries a marker exactly when some
summand does. Stable outputs are fixed by the cube endpoints. It therefore
suffices to show that a variable essential for any summand is essential
for $F_m$ on that cube.

If the lower corner already contains a full pair $\{a_h,b_h\}$,
nonconstancy forces it to contain no other variable: $x$ or any outside
bit fixed to $1$ would make the $h$-summand constantly $1$.
The restricted function is consequently the disjunction of $x$ and the
outside variables still free. Every free variable is essential.

Otherwise no pair is already full. Suppose a variable in pair $i$ is
essential for its summand. Set its mate to $1$ and leave all other
pairs incomplete at their lower values. Activate the second factor
by $x$ if available, otherwise by an outside bit already fixed to $1$,
otherwise by one available outside bit with its mate set to $0$.
The last choice is possible because in that case no outside bit was
fixed to $1$. Toggling the queried variable now changes $F_m$ from
$0$ to $1$.

If the essential variable is instead $x$ or an outside variable in the
second factor of summand $i$, no other input of that factor is fixed
to $1$. Set pair $i$ to $1$, all outside variables to $0$ except the
queried variable, and toggle that variable. A single outside bit cannot
complete another pair. Again $F_m$ changes from $0$ to $1$.
These pivotal edges prove canonicality.

The $x$-essential contexts are exactly the $m$ pure pair-sets
$P_i=\{a_i,b_i\}$: with $x=1$ some pair must be full, whereas with
$x=0$ a full pair wins as soon as any outside bit is present.
For $i\ne j$, the cube
\[
 L=\{a_j\},\qquad U=P_i\cup\{a_j\}
\]
contains no $P_h$, but $P_i\subseteq U$ and $L\subseteq P_j$.
The graph of \cref{core:prop:essential-context-coloring} is therefore
the complete graph on $m$ vertices. Its coloring bound gives $m$
occurrences; the displayed construction uses exactly $m$.
\end{proof}
